% ============================================================================
%  presentation/formulas.tex --- every equation in the deck, explained, EN | FR
%
%  COMPANION TO script.tex. That one tells you what to SAY; this one tells you
%  what you are saying when you point at an equation. Same two-column shape:
%  English on the left because the talk is in English, French on the right
%  because nobody should defend an equation they only half own.
%
%  EVERY NUMBER IN HERE WAS RECOMPUTED, not copied off a slide. Where a
%  recomputation disagreed with the deck it is flagged in the margin rather than
%  quietly reconciled -- there is one such case, the disc-overlap figure in
%  formula 9.
%
%  Build:  make -C presentation formulas    ->  formulas.pdf
% ============================================================================
\documentclass[10pt,a4paper]{article}

\usepackage{iftex}
\ifLuaTeX
  \usepackage{fontspec}
  \IfFontExistsTF{Lato}{\setsansfont{Lato}}{}
\else
  \usepackage[utf8]{inputenc}
  \usepackage[T1]{fontenc}
  \usepackage{lmodern}
\fi
\usepackage[english,french]{babel}
\usepackage[a4paper,margin=15mm,top=17mm,bottom=17mm]{geometry}
\usepackage{xcolor}
\usepackage{amsmath}
\usepackage{amssymb}
\usepackage{array}
\usepackage{booktabs}
\usepackage{microtype}
\usepackage{fancyhdr}
\usepackage[hidelinks]{hyperref}

\definecolor{ink}   {HTML}{1A1A1A}
\definecolor{accent}{HTML}{1F5C8B}
\definecolor{good}  {HTML}{2E7D32}
\definecolor{bad}   {HTML}{C62828}
\definecolor{warn}  {HTML}{EF6C00}
\definecolor{muted} {HTML}{6E6E6E}
\definecolor{rule}  {HTML}{D8D8D4}
\definecolor{wash}  {HTML}{F4F6F8}

\renewcommand{\familydefault}{\sfdefault}
\color{ink}
\setlength{\parindent}{0pt}
\setlength{\parskip}{0pt}

\pagestyle{fancy}
\fancyhf{}
\renewcommand{\headrulewidth}{0pt}
\fancyfoot[C]{\footnotesize\color{muted}\thepage}
\fancyhead[L]{\footnotesize\color{muted}Noise Modelling Hanoi --- the equations}
\fancyhead[R]{\footnotesize\color{muted}EN \; $\cdot$ \; FR}

% ---------------------------------------------------------------- structure
% \formula{n}{name}{where it appears}
\newcommand{\formula}[3]{%
  \par\vspace{5mm}\noindent
  \colorbox{accent}{\parbox{\dimexpr\textwidth-2\fboxsep}{%
    \color{white}\vspace{.6mm}
    \textbf{#1}\hspace{3mm}#2\hfill\small #3\vspace{.6mm}}}
  \par\vspace{2mm}}

\newcommand{\act}[2]{%
  \par\vspace{7mm}\noindent
  {\color{accent}\rule{\textwidth}{.9pt}}\par\vspace{1.5mm}
  {\large\bfseries\color{accent}#1}\hfill{\color{muted}#2}
  \par\vspace{.5mm}{\color{rule}\rule{\textwidth}{.4pt}}\par}

% The equation itself, on a wash ground so the eye finds it.
\newcommand{\eq}[1]{%
  \par\vspace{1mm}\noindent
  \colorbox{wash}{\parbox{\dimexpr\textwidth-2\fboxsep}{%
    \vspace{2mm}\begin{center}$\displaystyle #1$\end{center}\vspace{1mm}}}
  \par\vspace{2mm}}

% Symbol glossary. Two per row keeps it compact and still readable.
\newenvironment{where}{%
  \par\vspace{.5mm}{\footnotesize\color{muted}\textsc{where}}\par\vspace{.8mm}
  \begin{tabular}{@{}p{20mm}p{62mm}@{\hspace{6mm}}p{20mm}p{62mm}@{}}
}{%
  \end{tabular}\par\vspace{2mm}}

\newcolumntype{S}{>{\setlength{\parskip}{1.1mm}}p{84mm}}
\newenvironment{say}{%
  \begin{tabular}{@{}S@{\hspace{5mm}}|@{\hspace{5mm}}S@{}}
}{%
  \end{tabular}\par}

% The deck's shorthands, so an equation can be pasted here from main.tex and
% still compile. Same list as script.tex keeps.
\newcommand{\dB}{\ensuremath{\,\mathrm{dB}}}  % \ensuremath: text here, math in the deck
\newcommand{\rr}{$R^2$}
\newcommand{\LA}{$L_{A,25\mathrm{s}}$}

\newcommand{\must}[1]{{\color{good}\textbf{#1}}}
\newcommand{\never}[1]{{\color{bad}\textbf{#1}}}
\newcommand{\care}[1]{{\color{warn}\textbf{#1}}}

\newcommand{\flag}[2]{%
  \par\vspace{1.5mm}\noindent
  \colorbox{wash}{\parbox{\dimexpr\textwidth-2\fboxsep}{\vspace{1mm}
    \footnotesize\textbf{#1}\par\vspace{.7mm}#2\vspace{1mm}}}\par}

\begin{document}
\selectlanguage{english}

\begin{center}
  {\LARGE\bfseries\color{accent} The equations, explained}\\[1.5mm]
  {\large Every formula in the deck: what it says, why that form, and what it
   does not say}\\[3mm]
  {\color{muted}Companion to \texttt{script.pdf} \; $\cdot$ \; eleven formulas \;
   $\cdot$ \; English and French}
\end{center}

\vspace{4mm}
\noindent
\colorbox{wash}{\parbox{\dimexpr\textwidth-2\fboxsep}{\vspace{1.5mm}
\textbf{Why this exists.} \texttt{script.pdf} tells you what to \emph{say}. This
tells you what you are saying when you point at an equation --- and, more useful
under questioning, \textbf{what each one refuses to claim}. An acoustics audience
will not test whether you can recite a formula. They will test whether you know
which of its terms you actually fitted.

\medskip
\selectlanguage{french}
\textbf{À quoi ça sert.} \texttt{script.pdf} dit ce qu'il faut \emph{dire}. Ce
document dit ce que vous dites quand vous pointez une équation --- et, plus utile
en questions, \textbf{ce que chacune refuse d'affirmer}. Un public d'acousticiens
ne testera pas votre récitation. Il testera si vous savez lesquels de ses termes
vous avez réellement ajustés.
\selectlanguage{english}
\vspace{1.5mm}}}

\vspace{3mm}
\noindent
{\footnotesize\color{muted}Every number here was recomputed from the repository
rather than copied off a slide. One recomputation disagreed with the deck; it is
flagged in formula 9 rather than quietly reconciled.}

\clearpage
% ============================================================================
\act{Part I --- the quantity}{what a number in this study actually is}

\formula{1}{Sound pressure level}{the unit behind every figure in the talk}
\eq{L_p \;=\; 10\log_{10}\!\frac{p^{2}}{p_0^{2}}
    \;=\; 20\log_{10}\!\frac{p}{p_0},
    \qquad p_0 = 20\,\mu\mathrm{Pa}}
\begin{where}
  $p$ & sound pressure, in pascals &
  $p_0$ & the reference: the quietest pressure a healthy young ear detects at
          1\,kHz \\
\end{where}
\begin{say}
A decibel is not a unit of anything. It is \must{a ratio, on a log scale},
against a fixed reference. That is the whole definition.

Two consequences run through the rest of this document. Ten decibels is a factor
of ten \emph{in energy}. And because it is a ratio against a reference,
\must{a level is only meaningful if you know your reference} --- which, with an
uncalibrated microphone, we do not.

That is what ``relative, not absolute'' means. The \emph{shape} of our numbers is
sound. Their \emph{origin} on the scale is unknown to within a few decibels.
&
Un décibel n'est l'unité de rien. C'est \must{un rapport, en échelle
logarithmique}, contre une référence fixe. Toute la définition est là.

Deux conséquences traversent ce document. Dix décibels, c'est un facteur dix
\emph{en énergie}. Et comme c'est un rapport à une référence, \must{un niveau n'a
de sens que si l'on connaît sa référence} --- ce qui, avec un microphone non
étalonné, n'est pas notre cas.

C'est ça, «~relatif et non absolu~». La \emph{forme} de nos chiffres est bonne.
Leur \emph{origine} sur l'échelle nous est inconnue à quelques décibels près.
\end{say}

\formula{2}{The measured quantity, $L_{A,25\mathrm{s}}$}{slide 9 --- Metrology}
\eq{L_{A,25\mathrm{s}} \;=\; 10\log_{10}\!\left(\frac{1}{T}\int_{0}^{T}
    \frac{p_A^{2}(t)}{p_0^{2}}\,\mathrm{d}t\right),
    \qquad T \approx 25\,\mathrm{s}}
\begin{where}
  $p_A(t)$ & the pressure \emph{after} A-weighting (formula 3) &
  $T$ & the integration window: about 25 seconds \\
\end{where}
\begin{say}
An \emph{energy average over a window}, not an average of levels. The integral is
of $p^{2}$; the logarithm is taken once, at the end. Formula 4 says why that
order cannot be swapped.

\must{The window is the whole argument of slide 9.} Change $T$ and you have a
different quantity with a different name:

\quad $T = 25$\,s \; --- \; ours.\newline
\quad $T = 8$\,h \; --- \; $L_\mathrm{Aeq,8h}$, an occupational exposure.\newline
\quad $T = $ a year, with $+5$\,dB on the evening and $+10$\,dB at night \; ---
\; $L_\mathrm{den}$, what the WHO guidelines and the EU directive regulate.

\never{Comparing a 25 s sample to $L_\mathrm{den}$ is not a rounding error.} It is
comparing a photograph to a year of weather. That is why the WHO 53/45 values
were withdrawn from this work entirely.
&
Une \emph{moyenne énergétique sur une fenêtre}, pas une moyenne de niveaux.
L'intégrale porte sur $p^{2}$~; le logarithme est pris une seule fois, à la fin.
La formule 4 dit pourquoi cet ordre n'est pas interchangeable.

\must{La fenêtre est tout l'argument de la diapo 9.} Changez $T$ et vous avez une
autre grandeur, avec un autre nom~:

\quad $T = 25$\,s \; --- \; la nôtre.\newline
\quad $T = 8$\,h \; --- \; $L_\mathrm{Aeq,8h}$, une exposition professionnelle.\newline
\quad $T = $ un an, avec $+5$\,dB le soir et $+10$\,dB la nuit \; --- \;
$L_\mathrm{den}$, ce que réglementent l'OMS et la directive européenne.

\never{Comparer un échantillon de 25 s à $L_\mathrm{den}$ n'est pas une erreur
d'arrondi.} C'est comparer une photographie à une année de météo. D'où le retrait
complet des valeurs OMS 53/45 de ce travail.
\end{say}

\formula{3}{A-weighting, IEC 61672-1}{slide 45 --- backup}
\eq{R_A(f) = \frac{12194^{2}f^{4}}
      {(f^{2}+20.6^{2})\,\sqrt{(f^{2}+107.7^{2})(f^{2}+737.9^{2})}\,(f^{2}+12194^{2})},
    \qquad A(f) = 20\log_{10}R_A(f) + 2.00}
\begin{where}
  $f$ & frequency, in hertz &
  $+2.00$ & the offset that makes $A(1000\,\mathrm{Hz}) = 0$ exactly \\
\end{where}
\begin{say}
A fixed filter, standing in for the ear's unequal sensitivity across frequency.
Four poles and the normalising offset --- nothing is fitted, nothing is ours.

What it does: $-26$\,dB at 63\,Hz, $-16$\,dB at 125\,Hz, $0$ at 1\,kHz, $+1$\,dB
at 4\,kHz. \must{It discounts low frequencies heavily.}

That matters here more than it usually does. A motorcycle engine has far more
low-frequency content than a car, so A-weighting \care{discards a real part of
what makes a Hanoi street feel loud}. It is the right standard, it is what every
reference we compare against uses --- and it is a known bias of the quantity, not
of our protocol.

The app's implementation is pinned to the IEC table by a unit test, so the curve
is checked rather than assumed.
&
Un filtre fixe, qui tient lieu de la sensibilité inégale de l'oreille selon la
fréquence. Quatre pôles et l'offset de normalisation --- rien n'est ajusté, rien
n'est de nous.

Ce qu'il fait~: $-26$\,dB à 63\,Hz, $-16$\,dB à 125\,Hz, $0$ à 1\,kHz, $+1$\,dB à
4\,kHz. \must{Il déprécie fortement les basses fréquences.}

Ça compte ici plus qu'ailleurs. Un moteur de deux-roues a bien plus de contenu
basse fréquence qu'une voiture, donc la pondération A \care{écarte une part réelle
de ce qui rend une rue de Hanoï bruyante}. C'est la bonne norme, c'est celle de
toutes nos références --- et c'est un biais connu de la grandeur, pas de notre
protocole.

L'implémentation de l'application est verrouillée sur la table IEC par un test
unitaire~: la courbe est vérifiée, pas supposée.
\end{say}

\formula{4}{Decibels add in energy}{slide 45 --- backup}
\eq{L_{\Sigma} \;=\; 10\log_{10}\sum_{i} 10^{L_i/10},
    \qquad
    \bar L \;=\; 10\log_{10}\!\left(\frac{1}{n}\sum_{i} 10^{L_i/10}\right)}
\begin{say}
To combine levels you must leave the log scale, add, and come back.
\never{The arithmetic mean of decibels is not the mean level.}

It is not a small effect. Two sources of 70\,dB together give 73\,dB, not 70 and
not 140. And 60\,dB alongside 70\,dB gives 70.4 --- \must{the quiet one is
essentially invisible}, which is why a single horn dominates a 25 s window.

\must{Every arithmetic mistake this project made in the simulation came from
forgetting this line.} Formula 6 is one of them, made twice.
&
Pour combiner des niveaux il faut quitter l'échelle log, additionner, et revenir.
\never{La moyenne arithmétique de décibels n'est pas le niveau moyen.}

L'effet n'est pas petit. Deux sources de 70\,dB donnent ensemble 73\,dB, ni 70 ni
140. Et 60\,dB à côté de 70\,dB donne 70,4 --- \must{la source calme est
quasiment invisible}, ce qui explique qu'un seul klaxon domine une fenêtre de
25~s.

\must{Toutes les erreurs d'arithmétique de ce projet dans la simulation viennent
de l'oubli de cette ligne.} La formule 6 en est une, commise deux fois.
\end{say}

\clearpage
% ============================================================================
\act{Part II --- the model}{the three constants that beat everything else}

\formula{5}{The delivered kernel}{slide 16 --- The delivered model}
\eq{E \;=\; \frac{A_\mathrm{hw}}{\max(d_\mathrm{hw},D_0)}
        \;+\; \frac{A_\mathrm{res}}{\max(d_\mathrm{res},D_0)}
        \;+\; B,
    \qquad L \;=\; 10\log_{10}E}
\begin{where}
  $d_\mathrm{hw}$ & distance to the nearest highway-class road, in metres &
  $d_\mathrm{res}$ & distance to the nearest residential-class road \\
  $A_\mathrm{hw}$ & fitted, $4.774\times10^{7}$ &
  $A_\mathrm{res}$ & fitted, $3.800\times10^{7}$ \\
  $B$ & fitted background, $1.1\times10^{-10}$ &
  $D_0$ & 5\,m, a near-field cutoff, not fitted \\
\end{where}
\begin{say}
Two roads, three constants, and \must{no morphological variable at all}. Read it
in energy, not in decibels: each road contributes energy that falls off with
distance, they add (formula 4), and the logarithm is taken once at the end.

$\max(d,D_0)$ is not physics, it is arithmetic hygiene: at $d\to0$ the term
diverges. Five metres is roughly the width of a traffic lane --- closer than that,
``distance to the road'' has stopped meaning anything.

All three constants are constrained positive, so each stays interpretable as an
emitted energy.

\must{And say the honest part.} $B$ fits to $1.1\times10^{-10}$, which is
$-99.6$\,dB --- a boundary solution. \care{The background is not identifiable from
363 points over three sites.} A three-parameter model is, in practice, two.
&
Deux routes, trois constantes, et \must{aucune variable morphologique}. À lire en
énergie, pas en décibels~: chaque route apporte une énergie qui décroît avec la
distance, elles s'additionnent (formule 4), et le logarithme est pris une fois à
la fin.

$\max(d,D_0)$ n'est pas de la physique, c'est de l'hygiène arithmétique~: en
$d\to0$ le terme diverge. Cinq mètres, c'est à peu près la largeur d'une voie ---
plus près que ça, «~distance à la route~» ne veut plus rien dire.

Les trois constantes sont contraintes positives, donc chacune reste
interprétable comme une énergie émise.

\must{Et dites la partie honnête.} $B$ converge vers $1,1\times10^{-10}$, soit
$-99,6$\,dB --- une solution de bord. \care{Le fond n'est pas identifiable à partir
de 363 points sur trois sites.} Un modèle à trois paramètres en a, en pratique,
deux.
\end{say}

\formula{6}{Why $1/d$ and not $1/d^{2}$}{slide 16, and slide 45}
\eq{\text{point source: } E \propto \frac{1}{d^{2}} \;\Rightarrow\; -6.02\dB
    \text{ per doubling}
    \qquad
    \text{line source: } E \propto \frac{1}{d} \;\Rightarrow\; -3.01\dB
    \text{ per doubling}}
\begin{say}
A point source spreads its energy over a sphere, whose area goes as $d^{2}$. A
long line of sources spreads it over a cylinder, whose area goes as $d$.

\must{A street is a line, not a point.} A continuous stream of vehicles on a road
is much closer to an infinite line source than to one loudspeaker, so the energy
falls off as $1/d$ and the level drops about three decibels each time you double
your distance --- not six.

This is the one piece of real acoustics in the delivered model, and it is why a
two-parameter distance law does as well as it does. \must{It is not a curve we
fitted; it is the geometry we chose before fitting anything.}

If someone asks why not $1/d^{2}$: because you would be modelling a street as a
single vehicle.
&
Une source ponctuelle répartit son énergie sur une sphère, dont l'aire va en
$d^{2}$. Une longue ligne de sources la répartit sur un cylindre, dont l'aire va
en $d$.

\must{Une rue est une ligne, pas un point.} Un flux continu de véhicules est bien
plus proche d'une source linéique infinie que d'un haut-parleur, donc l'énergie
décroît en $1/d$ et le niveau perd environ trois décibels à chaque doublement de
distance --- pas six.

C'est le seul morceau de vraie acoustique du modèle livré, et c'est pourquoi une
loi de distance à deux paramètres s'en sort aussi bien. \must{Ce n'est pas une
courbe ajustée~; c'est la géométrie choisie avant tout ajustement.}

Si on demande pourquoi pas $1/d^{2}$~: parce que ce serait modéliser une rue comme
un seul véhicule.
\end{say}

\formula{7}{The traffic scenario --- and the mistake made twice}{slide 33}
\eq{\underbrace{L(k) = 10\log_{10}\!\big(k\,(E-E_\mathrm{res}) + E_\mathrm{res}\big)}_{\text{\color{good}correct}}
    \qquad\qquad
    \underbrace{L(k) = 10\log_{10}(kE) = L + 10\log_{10}k}_{\text{\color{bad}wrong}}}
\begin{where}
  $k$ & the traffic multiplier, $0.2$ to $3$ &
  $E$ & the cell's total energy at the measured traffic \\
  $E_\mathrm{res}$ & the zone's own residual ambience --- the 5th percentile of
    its levels &
  & \\
\end{where}
\begin{say}
Traffic is not all the noise there is. Split the cell's energy into a traffic
share and a residual ambience, \must{scale only the traffic share}, then take the
log.

The wrong version scales the total, so it scales the ambience too --- as if
removing cars also removed the air conditioners. It collapses to a constant
offset, $10\log_{10}k$, independent of place: at $k = 0.2$, exactly
$\mathbf{-6.99}$\,dB everywhere.

The correct version depends on how much of the energy was traffic. With the
residual at about \care{28\,\%} of total energy, $k=0.2$ gives $\mathbf{-3.7}$\,dB
--- \must{half the claimed benefit.}

\must{This error was made twice.} Found and corrected in Python in July; rebuilt
in Kotlin in August by people who knew about it. What caught it the second time
was a unit test asserting the invariant $L(1) = L$ --- which the wrong version
also satisfies, so the test that matters is the one at $k \neq 1$ against the
model itself. The two now agree to within $0.2$\,dB across the whole range.
&
Le trafic n'est pas tout le bruit. Séparez l'énergie de la cellule en une part
trafic et une ambiance résiduelle, \must{ne multipliez que la part trafic}, puis
prenez le log.

La version fausse multiplie le total, donc aussi l'ambiance --- comme si retirer
les voitures retirait les climatiseurs. Elle se réduit à un décalage constant,
$10\log_{10}k$, indépendant du lieu~: à $k = 0,2$, exactement
$\mathbf{-6,99}$\,dB partout.

La version correcte dépend de la part d'énergie qui était du trafic. Avec un
résiduel à environ \care{28~\%} de l'énergie totale, $k=0,2$ donne
$\mathbf{-3,7}$\,dB --- \must{la moitié du bénéfice annoncé.}

\must{Cette erreur a été commise deux fois.} Trouvée et corrigée en Python en
juillet~; reconstruite en Kotlin en août par des gens qui la connaissaient. Ce qui
l'a rattrapée la seconde fois, c'est un test sur l'invariant $L(1) = L$ --- que la
version fausse satisfait aussi, donc le test qui compte est celui à $k \neq 1$
contre le modèle lui-même. Les deux s'accordent maintenant à $0,2$\,dB près sur
toute la plage.
\end{say}

\clearpage
% ============================================================================
\act{Part III --- how it is evaluated}{the part the whole talk turns on}

\formula{8}{$R^2$, MAE, RMSE}{slides 15, 16, 46}
\eq{R^{2} = 1 - \frac{\sum_i (y_i-\hat y_i)^{2}}{\sum_i (y_i-\bar y)^{2}},
    \qquad
    \mathrm{MAE} = \frac{1}{n}\sum_i |y_i-\hat y_i|,
    \qquad
    \mathrm{RMSE} = \sqrt{\frac{1}{n}\sum_i (y_i-\hat y_i)^{2}}}
\begin{say}
$R^{2}$ compares your errors to the errors of the dumbest possible model:
\must{always predicting the mean}. One is perfect, zero is no better than the
mean.

And \must{it can go negative} --- which is why the deck has values like $-0.42$.
That is not a bug and not a bounded quantity. Because these are computed
\emph{out of fold}, a model can be actively worse than the mean on data it has
never seen. Every negative number in the model table means exactly that.

MAE is the one to quote in the room, because it is in decibels and people can
picture it. \must{5\,dB} --- a factor of three in energy.
&
$R^{2}$ compare vos erreurs à celles du modèle le plus bête possible~:
\must{prédire toujours la moyenne}. Un, c'est parfait~; zéro, ce n'est pas mieux
que la moyenne.

Et \must{il peut être négatif} --- d'où les valeurs comme $-0,42$ dans le deck. Ce
n'est ni un bug ni une grandeur bornée. Comme le calcul est fait \emph{hors pli},
un modèle peut être franchement pire que la moyenne sur des données jamais vues.
Chaque nombre négatif du tableau veut dire exactement ça.

La MAE est celle à citer en salle, parce qu'elle est en décibels et qu'on se la
représente. \must{5\,dB} --- un facteur trois en énergie.
\end{say}

\formula{9}{Buffered leave-one-out --- and why 300 m}{slide 12 --- the methodological core}
\eq{\text{for each point } i:\quad \text{train on } \{\,j : d(i,j) > r\,\},
    \qquad r = 300\,\mathrm{m} = \text{the feature aggregation radius}}
\begin{say}
The morphology features are computed over a disc of radius 300\,m around each
point. So two nearby points are partly built from \emph{the same buildings and
the same roads}. Put one in training and the other in test and the model is being
scored on near-copies of what it was taught.

\must{The buffer must equal the feature radius.} Not a rule of thumb --- an
identity. Anything smaller leaks by exactly the amount of disc the two points
share.

Two discs of radius $r$ whose centres are $s$ apart overlap by

\quad $\dfrac{A_\cap}{A} = \dfrac{2}{\pi}\left(\arccos\dfrac{s}{2r}
      - \dfrac{s}{2r}\sqrt{1-\left(\dfrac{s}{2r}\right)^{2}}\,\right)$

\must{This is the identity the retracted $R^2 = 0.45$ broke.} It grouped on 110\,m
cells while the features spanned 300\,m.
&
Les variables de morphologie sont calculées sur un disque de rayon 300\,m autour
de chaque point. Deux points proches sont donc construits en partie sur
\emph{les mêmes bâtiments et les mêmes routes}. Mettez l'un en entraînement et
l'autre en test~: le modèle est noté sur des quasi-copies de ce qu'on lui a
appris.

\must{Le tampon doit égaler le rayon des variables.} Pas une règle empirique ---
une identité. Tout ce qui est plus petit fuit d'exactement la part de disque que
les deux points partagent.

Deux disques de rayon $r$ dont les centres sont distants de $s$ se recouvrent de

\quad $\dfrac{A_\cap}{A} = \dfrac{2}{\pi}\left(\arccos\dfrac{s}{2r}
      - \dfrac{s}{2r}\sqrt{1-\left(\dfrac{s}{2r}\right)^{2}}\,\right)$

\must{C'est l'identité que le $R^2 = 0{,}45$ retiré a enfreinte.} Il regroupait sur
des cellules de 110\,m alors que les variables portaient sur 300\,m.
\end{say}
\flag{The overlap figures, recomputed --- and one correction to the deck}{%
With $r = 300$\,m: two points \textbf{40\,m} apart share \textbf{91.5\,\%} of the
disc, \textbf{70\,m} share \textbf{85.2\,\%}, \textbf{110\,m} share
\textbf{76.8\,\%}, \textbf{150\,m} share \textbf{68.5\,\%}.

\care{Slide 12 says two points 110\,m apart ``share more than 85\,\%''. The lens
formula gives 76.8\,\%; 85\,\% is reached at 71\,m.} The claim overstates the
overlap and should be corrected --- \must{the argument does not need it}, because
the real leak is worse than the average case: grouping on 110\,m cells puts points
\emph{either side of a cell boundary} in different folds, and those can be
centimetres apart, sharing essentially \textbf{100\,\%} of their disc. ---
\selectlanguage{french}\care{La diapo 12 dit «~plus de 85~\%~» à 110\,m~; la
formule donne 76,8~\%.} L'argument n'en a pas besoin~: la vraie fuite est pire que
le cas moyen, car deux points de part et d'autre d'une frontière de cellule
peuvent être à quelques centimètres et partager ~100~\%.\selectlanguage{english}}

\formula{10}{Bootstrapped by block, not by point}{slide 46 --- the intervals}
\eq{\text{resample the } 17 \text{ blocks with replacement},\;
    B = 2000\;\Rightarrow\;
    \text{CI} = \big[\,q_{2.5},\, q_{97.5}\,\big] \text{ of the recomputed statistic}}
\begin{say}
The interval is built by resampling \must{blocks}, not points.

Why it matters: points inside a block are spatially correlated, so they are not
independent draws. Resample points and the bootstrap thinks it has 363
independent observations when it has closer to 17. \care{The interval comes out
far too narrow, and the ranking looks far more certain than it is.}

This is what makes the honest version of slide 46 possible: the intervals overlap
almost entirely, and \must{``the physical kernel beats log-distance'' is a coin
flip} at this sample size. A by-point bootstrap would have hidden that.
&
L'intervalle est construit en rééchantillonnant des \must{blocs}, pas des points.

Pourquoi ça compte~: les points d'un même bloc sont spatialement corrélés, donc ce
ne sont pas des tirages indépendants. Rééchantillonnez des points et le bootstrap
croit disposer de 363 observations indépendantes alors qu'il en a plutôt 17.
\care{L'intervalle sort beaucoup trop étroit, et le classement paraît bien plus
certain qu'il ne l'est.}

C'est ce qui rend possible la version honnête de la diapo 46~: les intervalles se
recouvrent presque entièrement, et \must{«~le noyau physique bat la
log-distance~» est un tirage à pile ou face} à cette taille d'échantillon. Un
bootstrap par point l'aurait masqué.
\end{say}

\formula{11}{The ceiling on $R^{2}$}{slide 17 --- ``isn't 0.25 very low?''}
\eq{\Delta L \sim \mathcal{N}(0,\,2\sigma^{2})
    \;\Longrightarrow\;
    \sigma = \frac{\sqrt{\pi}}{2}\,\mathbb{E}|\Delta L| \approx 0.886\,\mathbb{E}|\Delta L|,
    \qquad
    R^{2}_{\max} = 1 - \frac{\sigma^{2}}{s^{2}}}
\begin{where}
  $\Delta L$ & the difference between two readings a few metres apart, at the
               same hour &
  $\sigma$ & the irreducible spread of a single reading \\
  $s^{2}$ & the total variance of the 363 measurements, $51.1\dB^{2}$ &
  & \\
\end{where}
\begin{say}
Two points a few metres apart at the same hour must be given \emph{the same}
prediction by any spatial model. So \must{whatever they disagree by is variance no
spatial model can ever explain}, and it caps $R^{2}$ from above.

The derivation is one line. If each reading is the true value plus independent
noise of sd $\sigma$, their difference has variance $2\sigma^{2}$; the mean
absolute value of a $\mathcal{N}(0,\tau^{2})$ is $\tau\sqrt{2/\pi}$; substitute
and rearrange.

\begin{center}\footnotesize
\begin{tabular}{@{}lrrrr@{}}
  \toprule
  pairs & $n$ & $\mathbb{E}|\Delta L|$ & $\sigma$ & $R^{2}_{\max}$ \\
  \midrule
  $<20$\,m & 29 & 4.23\dB & 3.75\dB & 0.72 \\
  $<40$\,m & 87 & 5.44\dB & 4.82\dB & \must{0.55} \\
  $<60$\,m & 229 & 6.09\dB & 5.40\dB & 0.43 \\
  \bottomrule
\end{tabular}
\end{center}

At the 40\,m scale the map works at, the ceiling is about \must{0.54}. The
delivered model reaches \must{0.246} --- \must{half of what is attainable, not a
tenth.}

\care{Say the caveat on the same breath.} Pairs at the same hour may be on
different days, so $\sigma$ carries day-to-day variation too and
\care{overestimates} pure measurement noise --- which makes the ceiling an
\emph{under}estimate. And 29 pairs is few. It is an order of magnitude, not a
bound.
&
Deux points à quelques mètres, à la même heure, doivent recevoir \emph{la même}
prédiction de tout modèle spatial. Donc \must{leur désaccord est une variance
qu'aucun modèle spatial ne pourra expliquer}, et il plafonne le $R^{2}$.

La dérivation tient en une ligne. Si chaque lecture est la vraie valeur plus un
bruit indépendant d'écart-type $\sigma$, leur différence a pour variance
$2\sigma^{2}$~; la valeur absolue moyenne d'une $\mathcal{N}(0,\tau^{2})$ vaut
$\tau\sqrt{2/\pi}$~; on substitue et on réarrange.

\medskip
Le tableau ci-contre donne les trois bandes de séparation. À l'échelle de 40\,m où
travaille la carte, le plafond est d'environ \must{0,54}. Le modèle livré atteint
\must{0,246} --- \must{la moitié de l'atteignable, pas le dixième.}

\care{Dites la réserve dans la même respiration.} Des paires à la même heure
peuvent être sur des jours différents, donc $\sigma$ porte aussi la variation jour
à jour et \care{surestime} le bruit de mesure pur --- ce qui fait du plafond une
\emph{sous}-estimation. Et 29 paires, c'est peu. C'est un ordre de grandeur, pas
une borne.
\end{say}

\clearpage
% ============================================================================
\act{If you only remember four lines}{}

\vspace{2mm}
\noindent
\begin{tabular}{@{}p{58mm}p{112mm}@{}}
  \toprule
  \must{Decibels add in energy.} &
    Leave the log scale, add, come back. Formula 4. Every arithmetic mistake this
    project made came from skipping it. \\
  \addlinespace
  \must{The buffer equals the feature radius.} &
    Formula 9. Not a heuristic --- an identity. Breaking it is what produced the
    retracted $R^{2} = 0.45$. \\
  \addlinespace
  \must{A street is a line, not a point.} &
    Formula 6. $-3$\,dB per doubling, not $-6$. The only real acoustics in the
    delivered model, and it was chosen before anything was fitted. \\
  \addlinespace
  \must{$R^{2}$ has a ceiling here, and it is about 0.54.} &
    Formula 11. Two points metres apart disagree by 4--5\,dB. The delivered 0.246
    is half of what is reachable. \\
  \bottomrule
\end{tabular}

\vspace{8mm}
\noindent
\colorbox{wash}{\parbox{\dimexpr\textwidth-2\fboxsep}{\vspace{1.5mm}
\textbf{Three things not to say about these equations}

\medskip
\begin{tabular}{@{}p{56mm}p{104mm}@{}}
  \never{``Our model is physically validated.''} &
    Formula 5 has one genuine physical choice --- the $1/d$ geometry. The three
    constants are fitted, and one of them ($B$) is not identifiable at all. Say
    \must{``a physically motivated form, with fitted constants''}. \\
  \addlinespace
  \never{``We measured $L_\mathrm{Aeq}$.''} &
    Formula 2. We measured $L_{A,25\mathrm{s}}$. If the window is not stated, the
    quantity is not stated. \\
  \addlinespace
  \never{``The ceiling proves 0.25 is good.''} &
    Formula 11 is an order of magnitude from 29 pairs, and it overestimates
    $\sigma$. It makes 0.25 \must{unsurprising}, not \must{good}. \\
\end{tabular}
\vspace{1.5mm}}}

\vspace{6mm}
\noindent
{\footnotesize\color{muted}Sources. Formulas 2--4: \texttt{docs/metrology.md} and
IEC 61672-1. Formula 5: \texttt{models/hybrid\_physical.json}, fitted by
\texttt{scripts/04\_evaluate\_models.py}. Formula 7:
\texttt{mobile/.../study/Scenario.kt}, invariant asserted in
\texttt{ScenarioTest.kt}. Formulas 8--10:
\texttt{scripts/04\_evaluate\_models.py} and \texttt{docs/methodology.md}.
Formula 11: computed in \texttt{scripts/12\_presentation\_figures.py} from
\texttt{data/processed/measurements.csv}.}

\end{document}
